% Chapter 16 -- Claim: Selling Options Earns a Reliable Premium
% Requires: amsmath, amssymb, booktabs, array, tabularx, longtable, graphicx, tcolorbox (breakable, skins), hyperref or url.
% Defines the callout box only if the preamble has not.
\providecolor{fenavy}{HTML}{1B2A6B}
\providecolor{feshade}{HTML}{F3F5FA}
\makeatletter
\@ifundefined{callout}{%
  \newtcolorbox{callout}{breakable, enhanced,
    colback=white, colframe=gray!70, boxrule=0.5pt, arc=1mm,
    borderline west={2.5pt}{0pt}{gray!70},
    left=8pt, right=8pt, top=4pt, bottom=4pt}%
}{}
\makeatother

\providecolor{fenavy}{HTML}{1B2A6B}
\providecolor{feshade}{HTML}{F3F5FA}
\makeatletter
% Shaded box for optional technical material that general readers may skip.
\@ifundefined{technote}{%
  \newtcolorbox{technote}{breakable, enhanced,
    colback=feshade, colframe=fenavy!60, boxrule=0.5pt, arc=1mm,
    left=8pt, right=8pt, top=4pt, bottom=4pt}%
}{}
\makeatother

\chapter{Claim: Selling Options Earns a Reliable Premium}
\label{ch:19}
\providecommand{\term}[1]{\emph{#1}}

Large institutions pay to protect themselves against a volatile stock market. Pension funds, endowments and other long-term investors buy put options\index{option!put} on stock indices so that a stock market crash does not wipe out years of gains. Insurance companies that sell annuities with guaranteed minimum returns buy options to hedge the promises they have made to their customers.

\begingroup
\small
\setlength{\LTpre}{6pt}\setlength{\LTpost}{6pt}
\begin{longtable}{@{}>{\raggedright\arraybackslash}p{0.08\linewidth}>{\raggedright\arraybackslash}p{\dimexpr(\linewidth-0.08\linewidth-4\tabcolsep)/2\relax}>{\raggedright\arraybackslash}p{\dimexpr(\linewidth-0.08\linewidth-4\tabcolsep)/2\relax}@{}}
\caption{Typical reasons to buy and to sell call and put options.}
\label{tab:18-buy-sell} \\
\toprule
 & Reasons to buy & Reasons to sell \\
\midrule
\endfirsthead
\multicolumn{3}{@{}l}{\emph{(continued)}} \\
\toprule
 & Reasons to buy & Reasons to sell \\
\midrule
\endhead
\midrule
\multicolumn{3}{r@{}}{\emph{(continued on the next page)}} \\
\endfoot
\bottomrule
\endlastfoot
Call &
Bet on a sharp rise at a small, known cost.\par\smallskip
Lock in a price for a purchase later.
&
Earn income on stocks already owned (a covered call).\par\smallskip
Expect the price to stay below the strike. \\*
\addlinespace
Put &
Insure a portfolio against a fall.\par\smallskip
Hedge a promise to limit losses, as pension funds and insurers do.
&
Collect an insurance premium (put-writing).\par\smallskip
Buy the stock only if it falls to the strike. \\*
\end{longtable}
\endgroup


A put option\index{option!put} is insurance against a market fall. Its seller collects a premium up front and pays out when the market drops. If the insurance is priced above the expected claims, the seller earns a steady profit, with an occasional very bad year. Carr and Wu (2009) measured the \textbf{variance risk premium}\index{variance risk premium}, the gap between the volatility priced into options and the volatility that follows, and found it large and persistent. Figure~\ref{fig:18-vrp} shows the same pattern from 1990 to 2025: in almost every year, the volatility priced into index options was higher than the volatility the market then delivered, but in 2008 and 2020 the market moved far more than option prices had allowed for.

% Figure: implied versus realized volatility, approximate annual averages.
% NOTE: the numbers below are APPROXIMATE annual averages (VIX, and realized
% volatility of the S&P 500 over the same calendar year), entered by hand to show
% the pattern. Replace them with the output of vrp_figure.py before publication.
\pgfplotstableread[row sep=\\]{
  year   vix   realized  gap \\
  1990   23.1   14.5    8.6 \\
  1991   18.4   14.3    4.1 \\
  1992   15.5    9.6    5.9 \\
  1993   12.7    8.6    4.1 \\
  1994   13.9    9.8    4.1 \\
  1995   12.4    7.8    4.6 \\
  1996   16.4   11.8    4.6 \\
  1997   22.4   18.0    4.4 \\
  1998   25.6   20.2    5.4 \\
  1999   24.4   18.0    6.4 \\
  2000   23.3   22.1    1.2 \\
  2001   25.8   21.6    4.2 \\
  2002   27.3   25.9    1.4 \\
  2003   22.0   17.0    5.0 \\
  2004   15.5   11.1    4.4 \\
  2005   12.8   10.3    2.5 \\
  2006   12.8   10.0    2.8 \\
  2007   17.5   16.0    1.5 \\
  2008   32.7   41.0   -8.3 \\
  2009   31.5   27.3    4.2 \\
  2010   22.5   18.0    4.5 \\
  2011   24.2   23.3    0.9 \\
  2012   17.8   12.7    5.1 \\
  2013   14.2   11.1    3.1 \\
  2014   14.2   11.4    2.8 \\
  2015   16.7   15.5    1.2 \\
  2016   15.8   13.1    2.7 \\
  2017   11.1    6.7    4.4 \\
  2018   16.6   17.1   -0.5 \\
  2019   15.4   12.5    2.9 \\
  2020   29.3   34.4   -5.1 \\
  2021   19.7   13.1    6.6 \\
  2022   25.6   24.2    1.4 \\
  2023   16.9   13.1    3.8 \\
  2024   15.5   12.7    2.8 \\
  2025   19.0   18.0    1.0 \\
}\vrpdata

\begin{figure}[htbp]
\centering
\scalebox{0.9}{%
\begin{tikzpicture}
\begin{groupplot}[
  group style={group size=1 by 2, vertical sep=1.6cm, x descriptions at=edge bottom},
  width=0.89\linewidth,
  xmin=1989.3, xmax=2025.7,
  xtick={1990,1995,...,2025},
  x tick label style={/pgf/number format/1000 sep={}, font=\scriptsize},
  y tick label style={font=\scriptsize},
  axis lines=left, grid=major, grid style={gray!20},
  title style={font=\small\bfseries, align=left},
  legend style={draw=none, fill=none, font=\footnotesize, cells={anchor=west}},
]
% ---------------- Panel A: implied versus realized ----------------
\nextgroupplot[height=5.6cm, ymin=0, ymax=52, ytick={0,10,20,30,40},
  ylabel={Volatility (\%)}, ylabel style={font=\footnotesize},
  title={A. Option prices usually expect more volatility than the market delivers},
  legend columns=2, legend style={at={(0.5,0.99)}, anchor=north, /tikz/every even column/.append style={column sep=8pt}}]
\addplot[fenavy, very thick, mark=*, mark size=1.1pt] table[x=year, y=vix] {\vrpdata};
\addlegendentry{implied volatility (VIX), yearly average}
\addplot[red!60!black, thick, mark=square*, mark size=1.0pt] table[x=year, y=realized] {\vrpdata};
\addlegendentry{realized volatility of the S\&P 500}
\node[font=\scriptsize, red!60!black, anchor=south] at (axis cs:2008,41.5) {2008};
\node[font=\scriptsize, red!60!black, anchor=south] at (axis cs:2020,35) {2020};
% ---------------- Panel B: the gap ----------------
\nextgroupplot[height=4.6cm, ymin=-12, ymax=16, ytick={-10,-5,0,5,10},
  ylabel={Implied $-$ realized (points)}, ylabel style={font=\footnotesize},
  xlabel={Year}, xlabel style={font=\footnotesize},
  title={B. A small gain in most years, a large loss in a few},
  ybar, bar width=4.2pt, bar shift=0pt,
  legend columns=3, legend style={at={(0.5,0.99)}, anchor=north, /tikz/every even column/.append style={column sep=8pt}}]
\addplot[fill=fenavy, draw=none] table[x=year, y expr={max(\thisrow{gap},0)}] {\vrpdata};
\addlegendentry{seller gains}
\addplot[fill=red!60!black, draw=none] table[x=year, y expr={min(\thisrow{gap},0)}] {\vrpdata};
\addlegendentry{seller loses}
\draw[gray, thin] (axis cs:1989.3,0) -- (axis cs:2025.7,0);
\addlegendimage{gray, dashed, line legend}
\addlegendentry{average, about 3 points}
\draw[gray, dashed] (axis cs:1989.3,3.0) -- (axis cs:2025.7,3.0);
\node[font=\scriptsize, red!60!black, anchor=north] at (axis cs:2008,-8.6) {2008};
\node[font=\scriptsize, red!60!black, anchor=north] at (axis cs:2020,-5.4) {2020};
\end{groupplot}
\end{tikzpicture}%
}
\caption{Implied volatility is usually higher than the volatility that follows. Panel A compares the yearly average of the VIX, the volatility priced into one-month S\&P 500 options, with the realized volatility of the S\&P 500 in the same year. Panel B shows the difference. A seller of index options gains in years when it is positive and loses when it is negative. Implied volatility was higher in all but a few years since 1990, by about 3 percentage points on average, but in 2008 and 2020 the market moved far more than option prices had allowed for. Numbers are approximate yearly averages.}
\label{fig:18-vrp}
\end{figure}


Option income can be a reliable reward for bearing a risk that others want to shed, or it can be small gains in most months paid for by large losses in a few, and the two can look alike in a backtest\index{backtest}. The evidence has to separate them, and the crashes of 2008, February 2018 and March 2020 show what the losses look like when they come. 


The claim on trial is: \textbf{put-writing (selling index puts) or writing covered calls earns more than the risk taken, because implied volatility is higher than the volatility that follows.}



The lab in Section~\ref{sec:16-lab} repeats Carr and Wu's (2009) measurement of the variance risk premium with the free VIX and daily U.S. market returns. It also rebuilds \term{put-writing} and covered-call strategies that can be compared with Cboe's published PUT and \term{BXM} indexes, to check whether selling options earns a reliable premium, and how much of that premium is simply stock-market risk.

By the end of the chapter you will be able to:

\begin{itemize}
  \item use put--call parity\index{put--call parity} to show that a covered call and a cash-secured put\index{option selling!cash-secured put} are the same position;
  \item define the VIX\index{VIX} and the variance risk premium\index{variance}\index{risk premium}, and estimate an option seller's expected gain\index{volatility!implied}\index{volatility!realized};
  \item judge an option strategy against a benchmark with the same market exposure, using up and down betas\index{beta}\index{beta!up and down} and a measure that cannot be gamed by selling crash insurance\index{manipulation-proof performance measure (MPPM)};
  \item judge what an option-income fund is likely to deliver, and what can go wrong\index{Black--Scholes--Merton model}\index{Greeks}.
\end{itemize}


\section{Foundations: the VIX and the option seller's payoff}
\label{sec:16-foundations}

Chapter~\ref{ch:18} introduced option payoffs, \term{put--call parity}, the Black--Scholes--Merton model, the Greeks and implied volatility (Sections~\ref{sec:14-bsm} to~\ref{sec:14-iv}). This chapter adds the VIX and the economics of selling options.
\subsection{The VIX}

The VIX,  plotted in Figure~\ref{fig:18-vrp}, measures the volatility priced into one-month S\&P 500 options. Its square is approximately the fair price of a 30-day \term{variance swap}\index{variance swap|textbf}, a contract that pays the gap between the variance that occurs and a level fixed in advance. The variance risk premium is the expected profit of the seller of such a swap.

\textbf{Worked example.} The index is at 100, the risk-free rate\index{risk-free rate} is 4\%, and a one-month at-the-money put trades at an implied volatility of 20\%, so its Black--Scholes price is \$2.14. Suppose volatility over the month turns out to be 16\%. At that volatility, the put would have been worth \$1.68, so the seller's expected gain is about $\$2.14 - \$1.68 = \$0.46$, roughly vega (\$0.115 per point of volatility) times the 4-point gap. That is only an \emph{expected} gain. If the index falls 10\% during the month, the put pays \$10, and the seller loses $\$10 - \$2.14 = \$7.86$, about seventeen months of expected gains. The shape of this bargain is the ``sell a put'' panel of Figure~\ref{fig:14-four-payoffs}.


\section{Why selling options pays, and how it is done}
\label{sec:16-how}

\textbf{Why insurance might be overpriced.} Many investors hold stock portfolios they cannot easily sell, and they want protection against crashes. Pension funds, insurers and issuers of structured products buy index puts. The dealers who sell them must hold capital against the risk and cannot fully hedge it, so this steady demand pushes option prices up, especially for out-of-the-money puts. Sellers are paid for bearing a risk that strikes at the worst times, when investors' wealth is low and losses hurt most. A premium for that risk is not a free lunch. It is compensation, and it should be large precisely because the losses come at bad times (Bollen and Whaley 2004; G\^arleanu, Pedersen and Poteshman 2009).


\textbf{The shape of the payoff.} Most ``option income'' products are covered calls or put-writing, and both keep much of the stock market's risk, with a beta of roughly 0.5 to 0.8, so their returns mix the reward for selling options with the stock market's own extra return, the equity premium\index{equity premium}. A covered call keeps the stock market's losses but caps its gains at the strike plus the premium (Figure~\ref{fig:16-cc}). In calm or gently rising months it beats the market. In strong rallies it lags. In crashes it loses almost as much as the market, softened only by the premium. Its returns are therefore \term{negatively skewed}\index{skewness|textbf}: most months bring small gains, and a few bring large losses.

\begin{figure}[htbp]
\centering
\scalebox{0.9}{%
\begin{tikzpicture}
\begin{axis}[
  width=0.61\linewidth, height=0.38\linewidth,
  axis lines=left, axis line style={-stealth},
  xmin=80, xmax=120, ymin=-20, ymax=21,
  xtick={80,85,...,120}, ytick={-20,-10,0,10,20},
  xlabel={Index at expiry}, ylabel={Profit per \$100 of index (\$)},
  clip=false,
]
  \addplot[gray, thin, domain=80:120, samples=2] {0};
  % Index alone
  \addplot[fegray, thick, dashed, domain=80:120, samples=2] {x - 100};
  % Covered call: index plus the $2.47 premium, capped at the strike
  \addplot[fenavy, very thick] coordinates {(80,-17.53) (100,2.47) (120,2.47)};
  \node[fegray, anchor=north west] at (axis cs:112,10) {index alone};
  \node[fenavy, anchor=north east, align=right] at (axis cs:120,1.5)
    {covered call:\\gain capped at \$2.47};
  \node[fenavy, anchor=north west, align=left] at (axis cs:81,20)
    {losses nearly as large\\as the index's};
\end{axis}
\end{tikzpicture}%
}
\caption{A covered call keeps almost all of the market's losses but caps its gains. A one-month at-the-money covered call against the index alone, with the index at 100, a 4\% interest rate and 20\% implied volatility, as in the worked example, at which the call is worth \$2.47. The premium caps the gain and cushions the loss by the same amount, so the covered call keeps almost all of the index's downside.}
\label{fig:16-cc}
\end{figure}

That shape creates a measurement problem. The Sharpe ratio uses only the average return and the standard deviation, so it cannot see skewness. Selling out-of-the-money options can raise a portfolio's Sharpe ratio while making its worst losses much larger. Section~\ref{sec:16-measure} sets out measures that are harder to game.


\section{The claim}
\label{sec:16-claim}

\begin{callout}
\textbf{Claim.} Implied volatility is higher than the volatility that follows, on average, so option sellers earn a reliable premium. Strategies that sell index options, such as covered calls and put-writing, deliver better risk-adjusted returns than holding the index.
\end{callout}

The claim has two parts, and each can succeed or fail on its own:

\begin{itemize}
  \item \textbf{The variance risk premium.} Implied variance is, on average, higher than the variance that follows. This is a statement about prices, and it can be tested directly by comparing the VIX with the volatility that follows.
  \item \textbf{The strategy premium.} Covered calls and put-writing earn more than their risk justifies. For this to hold, the variance premium must survive trading costs and practical details such as the timing of trades, and the strategies must be judged against an honest measure of their risk.
\end{itemize}

To test it, three points need to be settled:

\begin{itemize}
  \item \textbf{What counts as risk.} A strategy with a beta of 0.5 should be compared with a portfolio holding half the market and half Treasury bills, not with the whole market. Its negative skewness should be counted, not ignored.
  \item \textbf{What counts as reliable.} A premium that is positive in most months but wiped out in a few is reliable in one sense and not in another. The tests report both the share of months with a gain and the size of the worst losses.
  \item \textbf{What would count against it.} The variance risk premium fails if implied volatility is not significantly above realized volatility on average. The strategy claim fails if covered calls and put-writing earn no extra return after adjusting for beta, or if their advantage disappears once skewness is taken into account.
\end{itemize}

\section{The original evidence}
\label{sec:16-evidence}

\textbf{Option returns.} To see whether investors pay for volatility itself, Coval and Shumway (2001) bought S\&P 500 \textbf{straddles}, a call and a put with the same strike, arranged so that they did not depend on the market's direction. These positions lost about 3\% a week on average. Bakshi and Kapadia (2003) found the same for delta-hedged options: buyers lost money on average, because volatility rises when markets fall and investors pay to own it.

\textbf{The variance risk premium.} Carr and Wu (2009) computed the fair price of a variance swap from option prices, the method behind today's VIX, and found implied variance well above realized variance for the S\&P 500 and other indexes. Bollerslev, Tauchen and Zhou (2009) found that the size of this gap predicts stock returns over the next quarter, which suggests it is a price for economic uncertainty.

\textbf{Covered calls and put-writing.} Whaley (2002) designed Cboe's BuyWrite Index (BXM), which holds the S\&P 500 and sells a one-month call each month, and found that it earned returns close to the index's with much lower volatility. A later study for Cboe reported 11.8\% a year from 1988 to 2006 against 11.7\% for the S\&P 500, with a standard deviation of 9.3\% against 13.9\% (Callan Associates 2006). Results like these drew large sums into option-income products.

\textbf{The skeptics.} Broadie, Chernov and Johannes (2009) argued that large average losses on index puts do not by themselves show that puts are overpriced, because standard models with crashes predict very negative and very noisy put returns. Israelov and Nielsen (2015) split covered-call returns into stock-market exposure, a bet against volatility and an unintended bet on reversals, and found that the stock-market exposure drove most of the risk and return. A covered call, in other words, is mostly a stock-market position.

\section{Measuring option-strategy returns}
\label{sec:16-measure}

Option-selling strategies are the textbook case in which the usual performance measures mislead. Four corrections matter.

\textbf{Adjust for beta.} A covered call with a beta of 0.6 should earn about 60\% of the stock market's extra return even with no volatility premium. Compare it with a portfolio of 60\% in the market and 40\% in Treasury bills, or use its CAPM alpha (Chapter~\ref{ch:15}). Matching the market's return with less volatility may mean no extra return at all.

\textbf{Look at both betas.} Selling options makes market exposure lopsided. The \term{down beta}, measured in months when the market falls, is close to 1, so the strategy takes most of the loss. The \textbf{up beta}, measured in months when the market rises, is much lower, because the call caps the gains.

\textbf{Count the skewness.} Goetzmann, Ingersoll, Spiegel and Welch (2007) showed that a manager with no skill can raise a portfolio's Sharpe ratio by selling out-of-the-money options, because the Sharpe ratio ignores skewness. They proposed a \term{manipulation-proof performance measure (MPPM)}, which counts a large loss far more heavily than a large gain, so selling crash insurance cannot raise the score.

\begin{technote}
\textbf{Technical note: the MPPM.} The measure is
\begin{equation}
  \Theta = \frac{1}{(1 - \rho)\,\Delta t}\,\ln\!\left\{\frac{1}{T}\sum_{t=1}^{T}\left[\frac{1 + r_t}{1 + r_{f,t}}\right]^{1-\rho}\right\}
  \label{eq:16-mppm}
\end{equation}
where $r_t$ and $r_{f,t}$ are the returns of the strategy and of the risk-free asset over $T$ periods of length $\Delta t$ years. $\Theta$ is the safe extra return that an investor with risk aversion $\rho$ would accept in place of the strategy. With $\rho = 3$, a typical value, one large loss outweighs many small gains.
\end{technote}

\textbf{Use a long enough sample.} A strategy that loses heavily once a decade needs several decades of data (Section~\ref{sec:01-foundations}), and even twenty years of put returns are very noisy (Broadie, Chernov and Johannes 2009). A backtest from 2010 to 2019 would miss both of the big losses since 1990.


\section{Lab~\thechapter: Replicating Carr and Wu (2009) and the Cboe put-write and buy-write indices}
\label{sec:16-lab}

\textbf{Claim on trial:} selling index options earns a reliable premium. The lab replicates the variance risk premium of Carr and Wu (2009) and the put-write and buy-write strategies studied by Whaley (2002), using daily market returns from the Kenneth French library and the VIX from FRED, which starts in 1990. Pricing the options at the VIX overstates the premium, because the VIX includes the extra price of out-of-the-money puts, so the \texttt{-{}-iv-shift} option prices them a few points lower, and Cboe's PUT and BXM histories can be added for comparison.

\textbf{Part A, the variance risk premium.} At each month-end, the VIX against the market's realized volatility over the next 21 trading days, with a Newey--West $t$-statistic, the share of months with implied volatility above realized, and the worst months (Lab Figures~\ref{ch:19}.1 and~\ref{ch:19}.2).

\textbf{Part B, strategies.} An \textbf{at-the-money put-write} backed by Treasury bills, like Cboe's PUT, and a \textbf{covered call with a strike 2\% above the index}, compared with the market on return, volatility, Sharpe ratio, skewness, drawdown, beta and alpha (Lab Figure~\ref{ch:19}.3).

\textbf{Part C, robustness checks.} Returns in crises and in the market's five worst months, up and down betas, and the MPPM (Lab Figure~\ref{ch:19}.4).

\textbf{Part D, stocks versus volatility.} Each strategy's return split into the part explained by beta and the alpha that remains, compared with the market scaled to the same beta and with a pure short variance swap (Lab Figure~\ref{ch:19}.5).

\label{sec:16-stress}
\textbf{Robustness checks.} A few months decide the outcome: October 1987, August 1998, September 2008 to February 2009, and February and March 2020. In each, option sellers lost most of what the market lost, softened only by one month's premium of 1--3\%. Put-writing and covered calls often beat the market on the Sharpe ratio\index{Sharpe ratio}, but the advantage should shrink under the MPPM, and a beta of 0.6 alone can explain most of their return\index{skewness|textbf}. A strategy whose edge disappears after these adjustments is being paid for crash risk\index{crash risk} and stock-market risk, not earning alpha. The lab ends by asking you to reach your own verdict on the scale of Section~\ref{sec:01-method} and compare it with Section~\ref{sec:16-verdict}.

\section{Where the evidence stands}
\label{sec:16-verdict}

The published studies point to a clear pattern: \textbf{implied volatility has been higher than realized volatility on average, so the variance risk premium is real, but covered-call and put-writing strategies earn most of their return from stock-market exposure and lose most of what the market loses in crashes.} The table summarizes the evidence as of November 2026.

{\small
\begin{longtable}{@{}>{\raggedright\arraybackslash}p{0.19\linewidth}>{\raggedright\arraybackslash}p{0.13\linewidth}>{\raggedright\arraybackslash}p{\dimexpr 0.45\linewidth-6\tabcolsep\relax}>{\raggedright\arraybackslash}p{0.23\linewidth}@{}}
\toprule
Situation & Holds? & Why & Evidence \\
\midrule
\endfirsthead
\multicolumn{4}{@{}l}{\emph{(continued)}} \\
\toprule
Situation & Holds? & Why & Evidence \\
\midrule
\endhead
\midrule
\multicolumn{4}{r@{}}{\emph{(continued on the next page)}} \\
\endfoot
\bottomrule
\endlastfoot
Implied above realized variance & \textbf{Yes} & Implied variance has been well above realized variance on average for the S\&P 500 and other indexes. & Carr and Wu (2009) \\*
\addlinespace[3pt]
Option buyers pay for volatility & \textbf{Yes} & Zero-beta straddles lost about 3\% a week, and delta-hedged options lost money on average. & Coval and Shumway (2001); Bakshi and Kapadia (2003) \\
\addlinespace[3pt]
An economic reason for the premium & \textbf{Yes} & Demand for crash insurance raises option prices, and the premium predicts market returns. & Bollen and Whaley (2004); G\^arleanu, Pedersen and Poteshman (2009); Bollerslev, Tauchen and Zhou (2009) \\
\addlinespace[3pt]
Covered calls against the index & \textbf{Partly} & The BuyWrite Index matched the S\&P 500's return with lower volatility from 1988 to 2006. & Whaley (2002); Callan Associates (2006) \\
\addlinespace[3pt]
Covered calls after adjusting for beta & \textbf{Weakly} & Stock-market exposure drove most of the risk and return; the volatility bet was under 10\% of the risk. & Israelov and Nielsen (2015) \\
\addlinespace[3pt]
Risk adjustment that prices skewness & \textbf{Weakly} & Selling options can raise a Sharpe ratio without skill; measures that count crash losses show a smaller advantage. & Goetzmann, Ingersoll, Spiegel and Welch (2007) \\*
\addlinespace[3pt]
Mispricing, judged from put returns & \textbf{No} & Very negative, noisy put returns are expected in models with crashes, so they do not by themselves show mispricing. & Broadie, Chernov and Johannes (2009) \\
\end{longtable}
\addtocounter{table}{-1}% uncaptioned longtable: keep table numbers unchanged
}

\textbf{What would change the picture.} The case for option-income strategies would strengthen if their alpha against a benchmark with the same beta stayed significant after skewness and costs over a full cycle of crises. It would weaken if costs and fees consistently ate up the premium.

\label{sec:16-open}
\textbf{Open questions.} Options that expire on the day they are traded made up over 61\% of S\&P 500 index option volume in May 2025 (Cboe Global Markets 2025); have they changed the price of one-month volatility? And what do covered-call ETFs deliver after fees, spreads and taxes\index{taxes}, compared with an index position with the same beta?

\section{What practitioners and investors can use}
\label{sec:16-investors}

Three facts help interpret the high payouts that option-income funds advertise (Israelov and Nielsen 2014). Their payouts are not yield: the income comes with a cap on the fund's gains, so part of it is the market's own return, paid in a different shape. The fund is mostly a stock fund, with a beta of 0.5 to 0.7 and a down beta near 1, so compare it with an index-and-Treasury-bill mix with the same beta. And its worst months, such as 2008 and March 2020, say more than its Sharpe ratio. Ask a manager who sells volatility what the largest possible one-day loss is, and whether the strategy uses borrowed money\index{leverage} or margin that could be called in a volatility spike.

% Key terms of the chapter, in alphabetical order, for readers new to finance.
% A term defined in an earlier chapter points to the chapter, section and page.
\providecolor{FEMaroon}{RGB}{128,30,45}
\providecolor{FECream}{RGB}{253,248,238}
\providecolor{FEGold}{RGB}{196,150,60}
\makeatletter
\@ifundefined{keyterms}{%
  \newtcolorbox{keyterms}[1][Key terms in this chapter]{breakable, enhanced,
    colback=FECream, colframe=FEGold, boxrule=0.6pt, arc=2pt, leftrule=3pt,
    left=8pt, right=8pt, top=6pt, bottom=6pt, before skip=12pt, after skip=12pt,
    fonttitle=\bfseries, coltitle=white, colbacktitle=FEMaroon, title={#1}}%
}{}
\makeatother

\begin{keyterms}
\small
\begin{description}[leftmargin=0pt, itemsep=2pt, parsep=0pt, font=\normalfont\bfseries\color{FEMaroon}]
  \item[BXM.] The BXM is Cboe's index of a covered-call strategy that holds the S\&P~500 and sells a one-month call option on it each month.
  \item[Down beta and up beta.] Down beta is beta measured in months when the market falls, and up beta is beta measured in months when it rises.
  \item[Manipulation-proof performance measure (MPPM).] The MPPM is a performance measure that counts a large loss far more heavily than a large gain, so that selling options cannot inflate it.
  \item[Negatively skewed.] Returns are negatively skewed when most periods bring small gains and a few bring large losses.
  \item[Put--call parity.] See Chapter~\ref{ch:18}, Section~\ref{sec:14-parity}, page~\pageref{sec:14-parity}.
  \item[Put-writing.] Put-writing is a strategy that sells put options month after month and holds the strike in Treasury bills.
  \item[Variance swap.] A variance swap is a contract that pays the gap between the variance that occurs and a level fixed in advance.
\end{description}
\end{keyterms}
